% Replace the complete Introduction AND Related Work with this section.
% Keep later labels sec:3, sec:4, sec:5, sec:6 unchanged; numbering updates.
% 16 distinct citations here; method_and_evaluation_notes.tex supplies two more.
\section{Introduction}
\label{sec:1}

Business process compliance (BPC) concerns the conformance of business processes to applicable laws, regulations, standards, and organizational policies~\cite{hashmi2018we}. Design-time checking seeks to identify potential violations before execution. A central difficulty is connecting requirements written in natural language with the activities, responsibilities, and control flow represented in process models. Surveys of compliance approaches and rule languages describe a range of representations for this purpose~\cite{hashmi2018we,zasada2023evaluation}.

Formal approaches provide explicit representations for automated checking. Awad et al.~\cite{awad2008efficient} combine BPMN-Q with temporal logic, Elgammal et al.~\cite{elgammal2016formalizing} develop a pattern-based compliance language, and Governatori et al.~\cite{governatori2016semantic} investigate regulatory compliance checking using LegalRuleML. These approaches motivate a complementary question: how can the semantics required by a checker be obtained from regulatory text? Sleimi et al.~\cite{sleimi2018automated} study NLP-based extraction of semantic legal metadata, while Michel et al.~\cite{michel2022identification} investigate the identification of decision rules in legislative documents. Within design-time BPC, Winter et al.~\cite{winter2020assessing} assess process models against regulatory documents. Sun et al.~\cite{sun2024design} further combine multi-granularity regulatory semantics, including modality, actor, action, condition, constraint, and exception, with process-model information. Their representation provides the basis for the extraction task studied here.

Large language models (LLMs) offer another mechanism for interpreting regulatory text. De Kinderen and Winter~\cite{de2024towards} investigate prompt templates for legal GRL modeling. Barrientos et al. study regulatory-change impact analysis~\cite{barrientos2026impact} and the complexity of legal change for BPC~\cite{barrientos2026taming}. De Jure uses iterative LLM self-refinement for structured regulatory-rule extraction~\cite{guliani2026dejure}. LLMs have also been applied more directly to process-related tasks: Wang et al.~\cite{wang2025leveraging} generate process triple sequences from specifications for compliance risk identification; Park et al.~\cite{park2026extracting} extract declarative process constraints from natural-language descriptions; and K\"olbel et al.~\cite{koelbel2026context} combine BPMN models with external security knowledge for explanation and compliance assessment.

These studies establish relevant uses of LLMs, but their outputs and evaluation tasks differ. Our question is more specific: how does LLM-based semantic extraction behave when evaluated against a reconstructed non-LLM extractor under the same six-element BPC representation, reference annotations, and evaluation procedure, and how do the resulting records affect a fixed downstream checker? Separating extraction from checking makes both field-level errors and downstream limitations observable, without attributing changes in the checker to the extractor.

We investigate LLM-based Semantic Extraction (LLM-SE) within a three-part design-time BPC framework. Part~1 parses BPMN models into Process Records, while Part~2 independently transforms regulatory requirements into Rule Records. These are parallel upstream branches with no input dependency between them; Part~3 consumes both records for compliance checking. LLM-SE generates candidate records in a single call, followed by deterministic structural adaptation, evidence alignment, and validation. The records are then supplied to a shared checker for missing-action, incorrect-actor, and out-of-order violations. Our contributions are:
\begin{itemize}
    \item An implementation of LLM-SE for an existing fine-grained BPC representation, with source-linked evidence and an explicit deterministic postprocessing procedure.
    \item A controlled extraction comparison against a reconstructed Rules-Only baseline, together with a factorial analysis of semantic examples, semantic guidance, and additional output-discipline instructions.
    \item A downstream comparison using the same checking pipeline on a separate GDPR-based development benchmark, identifying both performance differences and remaining rule--process alignment limitations.
\end{itemize}

Section~\ref{sec:3} presents the problem formulation and framework. Section~\ref{sec:4} describes the methodology. Section~\ref{sec:5} reports the experiments and results, and Section~\ref{sec:6} concludes the paper.
